% !TeX root = ../main.tex

\chapter{Conclusion and Future Work}\label{chapter:conclusion}

\section{Summary}

This thesis investigated communication--computation overlap through a
benchmark-based evaluation on LRZ CoolMUC-4. Comb was extended with an
adjustable synthetic computation window between communication initiation
and completion, together with optional polling during computation. The
experiments compared Open MPI over TCP, MPICH over TCP, and Intel MPI over
verbs and TCP in a two-node, two-rank configuration across 19 peer payloads.
The four evaluated cases distinguish asynchronous progress settings and
final communication completion through \texttt{wait\_all} or
\texttt{test\_all}. Additional Intel MPI verbs experiments examined polling
during computation at four payloads.

The evaluation used changes in exposed receive completion time to quantify
measured overlap. Comb's mock backend provided a reference for local
framework and memory-handling costs. The resulting measurements characterize
how the tested communication configurations behave during an injected
computation window; they do not directly measure the execution time of a
complete ocean model.

\section{Conclusions}

The experimental results show that different MPI implementations, transport configurations, and progress settings exhibit clearly different communication--computation overlap behavior, and that these differences vary with message size. Here, measured overlap refers to the normalized receive-wait reduction $R_{\mathrm{wait}}$ defined in \autoref{sec:timing-metrics}. In addition, simply using nonblocking communication does not guarantee effective communication--computation overlap. Without asynchronous progress, Intel MPI verbs exhibits little sustained measured overlap across most of the payload range. Intel MPI TCP shows moderate but irregular overlap for small payloads, while the overlap is close to zero for most large payloads. Enabling asynchronous progress substantially increases $R_{\mathrm{wait}}$ for both transport configurations. For Intel MPI verbs, the receive-wait reduction is approximately equal to the measured computation time at many medium and large payloads; this alone does not establish that all communication is hidden.

The effect of asynchronous progress also depends on the MPI implementation and message size. For Open MPI TCP, enabling the TCP BTL progress thread has relatively little effect for small payloads, but preserves substantially more overlap at larger payloads than the configuration without the thread. Nevertheless, the measured overlap still decreases as the large payloads increase. MPICH shows an inconsistent response to asynchronous progress for small payloads, while the progress-enabled configuration maintains substantially higher overlap in the medium- and large-payload range. Thus, even with progress enabled, the tested MPI implementations do not exhibit a common message-size dependence.

The comparison of communication completion methods in \autoref{sec:completion-methods} shows that, when the injected computation time is approximately equal to the reference communication time, replacing \texttt{MPI\_Waitall} with a loop repeatedly calling \texttt{MPI\_Testall} does not provide a consistent advantage. The two methods generally produce similar measured overlap, although differences remain at individual payloads. Since both methods execute after the computation has finished, this result only reflects differences in the final communication completion path and does not imply that polling during computation is ineffective or that the two methods are statistically equivalent.

Polling within the computation window produces a different response. In the Intel MPI verbs experiments described in \autoref{sec:compute-polling}, when asynchronous progress is disabled, adding two polling points increases the run mean overlap at all four tested payloads. With asynchronous progress enabled, the additional benefit of polling is more limited and is not consistent across payloads.

This thesis experimentally characterizes the overlap of communication and
computation in the evaluated environment, providing
a basis for selecting and testing progress strategies for halo-exchange
workloads.

\section{Limitations and Future Work}

One major limitation of this study is the limited number of independent repetitions. Although each configuration was evaluated using multiple independent runs to observe run-to-run variability, the measurements for some message sizes, especially in the small-message region, exhibited substantial variability. Therefore, the current experiments are sufficient to reveal the main trends, but the statistical stability of small differences at individual message sizes and of highly variable data points remains limited. Future work could increase the number of independent runs and use confidence intervals or other statistical measures to determine more systematically whether the observed differences are stable.

A second limitation arises from the timing methodology itself. Timers were inserted around the communication initiation, computation window, and communication completion phases according to the execution structure of Comb, while the Mock backend was used to correct for local framework and memory-handling overhead as far as possible. However, these timing intervals do not directly correspond to the actual communication phases inside the MPI implementation. Internal MPI protocol processing, data progress, and some local processing may cross the timing boundaries defined in this study. Therefore, the measured communication times and overlap metrics should be interpreted as approximate measures based on program-observable timing intervals rather than as direct measurements of the internal network progress of MPI. Future work could combine the benchmark with lower-level tracing, MPI profiling, or hardware performance counters to further validate where and when communication actually occurs.

Third, the experimental parameters in this study were primarily chosen to construct a controlled benchmark environment and do not completely match the runtime parameters of a specific full ocean model. The message payload sizes, computation-window durations, computation-to-communication ratios, MPI-rank topology, and synthetic computation workload were deliberately controlled in order to study the effects of individual factors separately. Therefore, the overlap behavior observed in this study cannot be directly equated with application performance in a real ocean model. Future work should extract representative message sizes, computation times, and communication patterns from actual ocean-model workloads and validate the results using realistic computation kernels or complete applications.

In addition, the main experiments in this thesis are limited to a small set of MPI implementations, transport paths, and node configurations on LRZ CoolMUC-4. Extending the evaluation to more nodes, multiple MPI ranks per node, and different CPU and progress-thread resource allocations would allow the effects of shared-node resource contention and more complex communication topologies on communication progress to be investigated. The explicit polling experiments within the computation window are also currently limited to Intel MPI verbs, four payload sizes, and a fixed number of polling points. Future experiments could systematically vary the polling frequency and placement to determine under which conditions the additional communication progress provided by explicit polling is sufficient to offset the overhead of entering MPI during computation.

Finally, representative ocean-model kernels and application-level experiments are needed to verify whether the communication--computation overlap observed in Comb persists under realistic memory-access, threading, and numerical computation patterns, and whether this overlap can actually reduce the total execution time of a complete simulation.

